\section{Results}
\label{sec:results}

Placeholder results.
% value-check: {"tag": "peak-warming-median-range", "unit": "K", "lower": 1.65, "upper": 7.05}
Median peak warming across the scenarios ranges from around 1.7 to 7.0~K
(table~\ref{tab:peak-warming}, figure~\ref{fig:warming}).

\begin{figure}
\centering
\includegraphics[width=\textwidth]{@warming-figure@}
\caption{
    Projected warming relative to 1850--1900.
    a, b) Median (line) and 5th to 95th percentile range (shading)
    from 2000 to 2100 and from 2100 to 2500 (note the different y-axes).
    a) also shows the observed and human-induced warming assessed in
    Indicators of Global Climate Change 2025 \citep[IGCC;][]{forster_2026_igcc}.
    c, d) Year of peak warming (to 2500), on the same time axes as a) and b).
    e) Rate of warming (linear trend over ten years, as in IGCC)
    for each scenario over 2016--2025, 2026--2035, 2031--2040, 2036--2045 and 2041--2050.
    In 2016--2025, IGCC's assessed rate of human-induced warming is also shown
    (grey; white line: best estimate, box: \textit{likely} range).
    f--h) Peak warming (to 2500), warming in 2100 and warming in 2500.
    In the scenarios' boxes in c--h),
    the outer box shows the 5th to 95th percentile range,
    the inner boxes the 17th to 83rd and 33rd to 67th percentile ranges
    and the white line the median.
    In each of f--h), M and H are shown in their own panel, on their own scale.
    IGCC's values are global mean surface temperature,
    ours are global surface air temperature.
}
\label{fig:warming}
\end{figure}

\begin{table}
\caption{Peak warming (to 2500) relative to 1850--1900.}
\label{tab:peak-warming}
\centering
@peak-warming-table@
\end{table}
